% !TeX root = ../Book4.tex
\begin{figure}[htbp]
	\centering
	\begin{tikzpicture}[font=\small,scale=0.9]
		\begin{scope}[xshift=-3.5cm]
			\fill[GeometryBlueLight,even odd rule] (0,0) circle (1.4) (0,0) circle (0.55);
			\draw[GeometryGray] (0,0) circle (1.4);
			\draw[GeometryGray] (0,0) circle (0.55);
			\draw[GeometryBlue,thick,->] (0.55,0)--(1.4,0);
			\draw[GeometryBlue,dashed,thick,->] (45:0.95) arc (45:365:0.95);
			\node[below] at (0,-1.65) {环域 \(a<|z|<b\)};
			\node[above] at (0,1.65) {径向连接与核心闭路};
		\end{scope}
		\draw[->] (-1.6,0)--(0.1,0) node[midway,above] {\(w=\log(z/a)\)};
		\begin{scope}[xshift=0.7cm,yshift=-1.2cm]
			\fill[GeometryBlueLight] (0,0) rectangle (3.6,2.4);
			\draw[GeometryGray] (0,0) rectangle (3.6,2.4);
			\draw[GeometryBlue,thick,->] (0,1.05)--(3.6,1.05);
			\draw[GeometryBlue,dashed,thick,->] (2,0)--(2,2.4);
			\draw[->] (0.8,0)--(1.2,0);
			\draw[->] (0.8,2.4)--(1.2,2.4);
			\node[below] at (0,0) {\(0\)};
			\node[below] at (3.6,0) {\(L=\log(b/a)\)};
			\node[left] at (0,2.4) {\(2\pi\)};
			\node[above] at (1.8,2.75) {上下两边按平移识别};
			\node[below] at (1.8,-0.6) {连接族沿 \(s\)，核心族沿 \(t\)};
		\end{scope}
	\end{tikzpicture}
	\caption[环域的对数圆柱与两族曲线]{对数坐标使密度 \(|\mathrm{d}z|/|z|\) 成为平坦长度 \(|\mathrm{d}w|\)。实线表示连接两个边界的曲线，虚线表示核心闭路；右图的上下边识别后，虚线的两端接成一圈。图示尺寸只表达坐标方向，模量由 \(L/(2\pi)\) 给出。}
	\label{g4:fig:0505:annulus-cylinder}
\end{figure}
